\documentclass{article}
\begin{document}
\title{Radio reconfiguration can wait for the evidence it would remove}
\maketitle

\paragraph{Selected source and cut.}
Use the pinned SubT visibility-route and broker fixture on five
sampled cave tiles. The selected range is 72 m, not the stock
parameter. No delay, map blockage, link fade, robot mobility or
range guarantee is measured. There are two owed two-byte reports.
The first delivery window is its specified tick \(p\) and the
second is tick \(b>p\). Near-site registration changes the route
both reports use. It selects a cheaper route for the old report
whose longest hop exceeds the scoped range; the new report becomes
deliverable. The stock broker and receiving callback therefore
deliver both exactly when an effective near-site registration at
\(k\) meets \(p<k\le b\), with activation preceding any dispatch
at the same tick. A separate far-site option does not break the
old route. The claim about the near-site cut is for this selected
source/geometry/links and sufficient bandwidth only.

\paragraph{Causal placement and observation.}
Let the authorized caller command at tick \(s\) with integer
activation delay \(\lambda\in[L,U]\), \(1\le L\le U\). The
old receiver gets an authentic application callback if the
first dispatch actually succeeds. A caller with routed access
learns this at \(r\ge p\); the location/forwarding delay is
explicitly modeled, not extracted from native broker timing.
A receipt-first command is robust exactly when \(r+U\le b\).
It cannot violate \(p<k\) in this envelope. Conversely the
latest delay and earliest permitted command give failure if
that inequality is false. This is a receipt-first bound, not
a general impossibility for all possible foreknowledge or
independent alternatives. For a caller without a qualifying
message or other separating evidence, one fixed cut must lie
in the intersection of the safe service windows of the possible
histories.

For example let the two potential report-window pairs be \((1,2)\)
and \((3,4)\), with equal initial local views and unit activation
latency. The source histories require activation at 2 and 4
respectively. An authorized caller co-located with the old
receiver can command on receiving \(P\) in both cases. A remote
caller that lacks a distinguishing observation until one tick
later has identical information at the last command cut in the
early case. Acting by that cut breaks the late case; not acting
loses the early report \(Q\). A fixed valid far site avoids
this particular obstruction if separately available, which
the replay checks too. The experiment does not claim the
stock runtime provides the staged activation authority on
those deadlines.

\paragraph{Relation to submission-generated source loss.}
The prior ARIAC run observed a valid pose before submitting
an order that both starts a new duty and suppresses its only
producer. The cave fixture instead waits for actual settlement
of an old radio report before installing a route that prevents
that settlement and opens a new one. The shared guard is to
retain authorized evidence of the last effect needed from the
old context, revalidate its prospective use where applicable,
then put the context change before the new opportunity expires.
Here the old callback clears its named delivery rather than
being a fresh geometry read; there is no general pose-to-receipt
substitution. The source-bound interpretation and the
deadline/delegation conditions are essential in each. The
same-information direct route planner ties. This is a selected
development case, not a frozen independently selected transfer
or a new mission acceptance result.
\end{document}
